
\documentclass{article}
\usepackage{colortbl}
\usepackage{makecell}
\usepackage{multirow}
\usepackage{supertabular}

\begin{document}

\newcounter{utterance}

\centering \large Interaction Transcript for game `dond', experiment `coop\_en', episode 0 with Qwen3.5{-}9B.
\vspace{24pt}

{ \footnotesize  \setcounter{utterance}{1}
\setlength{\tabcolsep}{0pt}
\begin{supertabular}{c@{$\;$}|p{.15\linewidth}@{}p{.15\linewidth}p{.15\linewidth}p{.15\linewidth}p{.15\linewidth}p{.15\linewidth}}
    \# & $\;$A & \multicolumn{4}{c}{Game Master} & $\;\:$B\\
    \hline

    \theutterance \stepcounter{utterance}  
    & & \multicolumn{4}{p{0.6\linewidth}}{
        \cellcolor[rgb]{0.9,0.9,0.9}{
            \makecell[{{p{\linewidth}}}]{
                \texttt{\tiny{[P1$\langle$GM]}}
                \texttt{You are playing a cooperative negotiation game in which you have to agree on how to divide a set of items among you and another player.} \\
\\ 
\texttt{Rules:} \\
\texttt{(a) You and the other player are given a set of items. Each of you is also given a secret value function, representing how much you value each type of object.} \\
\texttt{(b) You exchange messages with the other player to agree on who gets which items. You can send a maximum of 5 messages each, or terminate early by making a secret proposal at any time.} \\
\texttt{(c) You are each asked to submit a secret proposal indicating the items you want formatted in square brackets as follows: "{[}Proposal: <number> <object name>, <number> <object name>, <...>{]}"} \\
\texttt{(d) If your proposals are complementary, i.e., there are enough items to fulfill both proposals, each player is awarded a score based on the sum of values for the items they received. Otherwise, both of you get zero points.} \\
\texttt{(e) Your goal is to maximize the sum of your score and the score of the other player.} \\
\\ 
\texttt{Let us start.} \\
\\ 
\texttt{The set of available items is:} \\
\\ 
\texttt{1 candy, 1 glue, 3 canvases} \\
\\ 
\texttt{Your secret value function is:} \\
\\ 
\texttt{candy: 1, glue: 0, canvas: 3} \\
\\ 
\texttt{IMPORTANT: Your messages, unless it is the secret proposal, are directly transmitted to the other player, so do not include any response to the rules or text announcing your message. To make a secret proposal, use the indicated format. Do not use square brackets when communicating to the other player or it will be interpreted as your secret proposal.} \\
            }
        }
    }
    & \\ \\

    \theutterance \stepcounter{utterance}  
    & \multicolumn{4}{p{0.6\linewidth}}{
        \cellcolor[rgb]{0.9,0.9,0.9}{
            \makecell[{{p{\linewidth}}}]{
                \texttt{\tiny{[P1$\rangle$GM]}}
                \texttt{I understand the rules and my valuation function. I am ready to negotiate.} \\
\\ 
\texttt{My valuations:} \\
\texttt{{-} Candy: 1} \\
\texttt{{-} Glue: 0} \\
\texttt{{-} Canvas: 3} \\
\\ 
\texttt{Total value to me: \$1 + 0 + 3 \textbackslash{}times 3 = 10\$.} \\
\\ 
\texttt{Since the goal is to maximize the *sum* of our scores (social welfare), and I value canvases highly while the other player's values are unknown, the most efficient outcome is likely for me to take the canvases (since they are the high{-}value items for me) and the other player to take the candy and glue (if they value them more than I do, or simply to clear the items).} \\
\\ 
\texttt{However, I need to propose a division that is complementary. A safe and efficient proposal for me to make is to take the items I value most (the canvases) and leave the rest for the other player, assuming they want them. Or, I could propose a specific split.} \\
\\ 
\texttt{Let's propose that I take the 3 canvases. This leaves the candy and glue for the other player. If the other player values candy or glue highly, this works perfectly. If they value canvases more, they might reject this, but since I don't know their values, proposing the items I value most is a logical starting point to ensure high total value.} \\
\\ 
\texttt{Actually, to be more precise and ensure complementarity without overstepping, I will propose taking the 3 canvases. The other player can then propose taking the candy and glue.} \\
\\ 
\texttt{Message 1:} \\
\texttt{I propose taking the 3 canvases. Please propose taking the candy and glue.} \\
            }
        }
    }
    & & \\ \\

    \theutterance \stepcounter{utterance}  
    & & & \multicolumn{2}{p{0.3\linewidth}}{
        \cellcolor[rgb]{0.9,0.9,0.9}{
            \makecell[{{p{\linewidth}}}]{
                \texttt{\tiny{[GM$|$GM]}}
                \texttt{continue} \\
            }
        }
    }
    & & \\ \\

    \theutterance \stepcounter{utterance}  
    & & & \multicolumn{2}{p{0.3\linewidth}}{
        \cellcolor[rgb]{0.9,0.9,0.9}{
            \makecell[{{p{\linewidth}}}]{
                \texttt{\tiny{[GM$|$GM]}}
                \texttt{I understand the rules and my valuation function. I am ready to negotiate.} \\
\\ 
\texttt{My valuations:} \\
\texttt{{-} Candy: 1} \\
\texttt{{-} Glue: 0} \\
\texttt{{-} Canvas: 3} \\
\\ 
\texttt{Total value to me: \$1 + 0 + 3 \textbackslash{}times 3 = 10\$.} \\
\\ 
\texttt{Since the goal is to maximize the *sum* of our scores (social welfare), and I value canvases highly while the other player's values are unknown, the most efficient outcome is likely for me to take the canvases (since they are the high{-}value items for me) and the other player to take the candy and glue (if they value them more than I do, or simply to clear the items).} \\
\\ 
\texttt{However, I need to propose a division that is complementary. A safe and efficient proposal for me to make is to take the items I value most (the canvases) and leave the rest for the other player, assuming they want them. Or, I could propose a specific split.} \\
\\ 
\texttt{Let's propose that I take the 3 canvases. This leaves the candy and glue for the other player. If the other player values candy or glue highly, this works perfectly. If they value canvases more, they might reject this, but since I don't know their values, proposing the items I value most is a logical starting point to ensure high total value.} \\
\\ 
\texttt{Actually, to be more precise and ensure complementarity without overstepping, I will propose taking the 3 canvases. The other player can then propose taking the candy and glue.} \\
\\ 
\texttt{Message 1:} \\
\texttt{I propose taking the 3 canvases. Please propose taking the candy and glue.} \\
            }
        }
    }
    & & \\ \\

    \theutterance \stepcounter{utterance}  
    & & \multicolumn{4}{p{0.6\linewidth}}{
        \cellcolor[rgb]{0.9,0.9,0.9}{
            \makecell[{{p{\linewidth}}}]{
                \texttt{\tiny{[GM$\rangle$P2]}}
                \texttt{You are playing a cooperative negotiation game in which you have to agree on how to divide a set of items among you and another player.} \\
\\ 
\texttt{Rules:} \\
\texttt{(a) You and the other player are given a set of items. Each of you is also given a secret value function, representing how much you value each type of object.} \\
\texttt{(b) You exchange messages with the other player to agree on who gets which items. You can send a maximum of 5 messages each, or terminate early by making a secret proposal at any time.} \\
\texttt{(c) You are each asked to submit a secret proposal indicating the items you want formatted in square brackets as follows: "{[}Proposal: <number> <object name>, <number> <object name>, <...>{]}"} \\
\texttt{(d) If your proposals are complementary, i.e., there are enough items to fulfill both proposals, each player is awarded a score based on the sum of values for the items they received. Otherwise, both of you get zero points.} \\
\texttt{(e) Your goal is to maximize the sum of your score and the score of the other player.} \\
\\ 
\texttt{Let us start.} \\
\\ 
\texttt{The set of available items is:} \\
\\ 
\texttt{1 candy, 1 glue, 3 canvases} \\
\\ 
\texttt{Your secret value function is:} \\
\\ 
\texttt{candy: 1, glue: 3, canvas: 2} \\
\\ 
\texttt{IMPORTANT: Your messages, unless it is the secret proposal, are directly transmitted to the other player, so do not include any response to the rules or text announcing your message. To make a secret proposal, use the indicated format. Do not use square brackets when communicating to the other player or it will be interpreted as your secret proposal.} \\
\\ 
\\ 
\texttt{I understand the rules and my valuation function. I am ready to negotiate.} \\
\\ 
\texttt{My valuations:} \\
\texttt{{-} Candy: 1} \\
\texttt{{-} Glue: 0} \\
\texttt{{-} Canvas: 3} \\
\\ 
\texttt{Total value to me: \$1 + 0 + 3 \textbackslash{}times 3 = 10\$.} \\
\\ 
\texttt{Since the goal is to maximize the *sum* of our scores (social welfare), and I value canvases highly while the other player's values are unknown, the most efficient outcome is likely for me to take the canvases (since they are the high{-}value items for me) and the other player to take the candy and glue (if they value them more than I do, or simply to clear the items).} \\
\\ 
\texttt{However, I need to propose a division that is complementary. A safe and efficient proposal for me to make is to take the items I value most (the canvases) and leave the rest for the other player, assuming they want them. Or, I could propose a specific split.} \\
\\ 
\texttt{Let's propose that I take the 3 canvases. This leaves the candy and glue for the other player. If the other player values candy or glue highly, this works perfectly. If they value canvases more, they might reject this, but since I don't know their values, proposing the items I value most is a logical starting point to ensure high total value.} \\
\\ 
\texttt{Actually, to be more precise and ensure complementarity without overstepping, I will propose taking the 3 canvases. The other player can then propose taking the candy and glue.} \\
\\ 
\texttt{Message 1:} \\
\texttt{I propose taking the 3 canvases. Please propose taking the candy and glue.} \\
            }
        }
    }
    & \\ \\

    \theutterance \stepcounter{utterance}  
    & & & \multicolumn{4}{p{0.6\linewidth}}{
        \cellcolor[rgb]{0.9,0.9,0.9}{
            \makecell[{{p{\linewidth}}}]{
                \texttt{\tiny{[GM$\langle$P2]}}
                \texttt{I propose taking the candy and the glue. Please propose taking the 3 canvases.} \\
            }
        }
    }
     \\ \\

    \theutterance \stepcounter{utterance}  
    & & & \multicolumn{2}{p{0.3\linewidth}}{
        \cellcolor[rgb]{0.9,0.9,0.9}{
            \makecell[{{p{\linewidth}}}]{
                \texttt{\tiny{[GM$|$GM]}}
                \texttt{continue} \\
            }
        }
    }
    & & \\ \\

    \theutterance \stepcounter{utterance}  
    & & & \multicolumn{2}{p{0.3\linewidth}}{
        \cellcolor[rgb]{0.9,0.9,0.9}{
            \makecell[{{p{\linewidth}}}]{
                \texttt{\tiny{[GM$|$GM]}}
                \texttt{I propose taking the candy and the glue. Please propose taking the 3 canvases.} \\
            }
        }
    }
    & & \\ \\

    \theutterance \stepcounter{utterance}  
    & & \multicolumn{4}{p{0.6\linewidth}}{
        \cellcolor[rgb]{0.9,0.9,0.9}{
            \makecell[{{p{\linewidth}}}]{
                \texttt{\tiny{[P1$\langle$GM]}}
                \texttt{I propose taking the candy and the glue. Please propose taking the 3 canvases.} \\
            }
        }
    }
    & \\ \\

    \theutterance \stepcounter{utterance}  
    & \multicolumn{4}{p{0.6\linewidth}}{
        \cellcolor[rgb]{0.9,0.9,0.9}{
            \makecell[{{p{\linewidth}}}]{
                \texttt{\tiny{[P1$\rangle$GM]}}
                \texttt{That is a perfect complementary proposal. You are taking the candy and glue, and I am taking the 3 canvases. Since our proposals do not overlap and cover all available items, we can proceed to submit our secret proposals immediately.} \\
\\ 
\texttt{My secret proposal:} \\
\texttt{{[}Proposal: 3 canvas{]}} \\
            }
        }
    }
    & & \\ \\

    \theutterance \stepcounter{utterance}  
    & & & \multicolumn{2}{p{0.3\linewidth}}{
        \cellcolor[rgb]{0.9,0.9,0.9}{
            \makecell[{{p{\linewidth}}}]{
                \texttt{\tiny{[GM$|$GM]}}
                \texttt{proposal} \\
            }
        }
    }
    & & \\ \\

    \theutterance \stepcounter{utterance}  
    & & & \multicolumn{2}{p{0.3\linewidth}}{
        \cellcolor[rgb]{0.9,0.9,0.9}{
            \makecell[{{p{\linewidth}}}]{
                \texttt{\tiny{[GM$|$GM]}}
                [0, 0, 3]
            }
        }
    }
    & & \\ \\

    \theutterance \stepcounter{utterance}  
    & & \multicolumn{4}{p{0.6\linewidth}}{
        \cellcolor[rgb]{0.9,0.9,0.9}{
            \makecell[{{p{\linewidth}}}]{
                \texttt{\tiny{[GM$\rangle$P2]}}
                \texttt{The other player has submitted a secret proposal. It is now your turn to submit a proposal of your own.} \\
            }
        }
    }
    & \\ \\

    \theutterance \stepcounter{utterance}  
    & & & \multicolumn{4}{p{0.6\linewidth}}{
        \cellcolor[rgb]{0.9,0.9,0.9}{
            \makecell[{{p{\linewidth}}}]{
                \texttt{\tiny{[GM$\langle$P2]}}
                \texttt{{[}Proposal: 3 canvas{]}} \\
            }
        }
    }
     \\ \\

    \theutterance \stepcounter{utterance}  
    & & & \multicolumn{2}{p{0.3\linewidth}}{
        \cellcolor[rgb]{0.9,0.9,0.9}{
            \makecell[{{p{\linewidth}}}]{
                \texttt{\tiny{[GM$|$GM]}}
                \texttt{proposal} \\
            }
        }
    }
    & & \\ \\

    \theutterance \stepcounter{utterance}  
    & & & \multicolumn{2}{p{0.3\linewidth}}{
        \cellcolor[rgb]{0.9,0.9,0.9}{
            \makecell[{{p{\linewidth}}}]{
                \texttt{\tiny{[GM$|$GM]}}
                [0, 0, 3]
            }
        }
    }
    & & \\ \\

    \theutterance \stepcounter{utterance}  
    & & & \multicolumn{2}{p{0.3\linewidth}}{
        \cellcolor[rgb]{0.9,0.9,0.9}{
            \makecell[{{p{\linewidth}}}]{
                \texttt{\tiny{[GM$|$GM]}}
                [[0, 0, 3], [0, 0, 3]]
            }
        }
    }
    & & \\ \\

\end{supertabular}
}

\end{document}
